\chapter{Review of Related Literature}

This chapter reviews relevant literature on note-taking strategies, memory processes, natural language processing (NLP) for semantic similarity, and progress tracking techniques. These areas provide the theoretical basis for understanding information encoding, memory consolidation, textual comparison, and user engagement. By examining these interconnected concepts, the review establishes a foundation for exploring methods that enhance learning, recall, and performance. The studies discussed offer insights into optimizing knowledge retention, evaluating recall accuracy, and motivating users through structured feedback mechanisms.

\section{Memory and Cognitive Processes}

Memory plays a crucial role in the learning process, specifically in the retention and recall of information. Studies have demonstrated that memory prioritizes the retention of meaning over surface-level details, which has direct implications for comprehension and learning.

{\color{maroon}
\subsection{Working Memory and Learning Efficiency}

Klemm (2007) establishes that effective thinking does not occur in isolation but relies fundamentally on working memory capacity. Working memory, defined as small amounts of information accessible over short periods that individuals use in learning and thinking, serves as the cognitive foundation for intellectual performance. This "scratch pad" memory operates through successive operations where chunks of information are processed and integrated into the thinking process.

The correlation between working memory capacity and intellectual capabilities demonstrates the critical importance of memory in educational contexts. Research reveals tight correlations between IQ and working memory capacity, and between working memory and problem-solving abilities. These findings suggest that students' learning effectiveness depends not merely on exposure to information, but on their ability to maintain and manipulate that information within working memory during cognitive tasks.

Klemm's analysis challenges contemporary educational approaches that emphasize understanding while neglecting memorization skills. The prevailing focus on insight, creativity, and inquiry learning, while valuable, creates an educational bias that underappreciates memory's role in learning. Students cannot effectively apply concepts they understand if they cannot remember them when needed. Moreover, robust memory expands the repertoire of cognitive capabilities upon which new understandings can be developed and expedited.
}

\subsection{The Role of Semantic and Syntactic Memory}

Sachs (1967) explored the relationship between semantic and syntactic memory in the context of comprehension. In her study, participants were presented with sentences and tested on their ability to recognize changes in meaning and form after brief intervals. The findings revealed that memory retention is primarily focused on semantic content, while syntactic details, such as active-to-passive transformations, fade quickly. Even when participants retained the overall meaning of the sentences, the formal properties of the sentences were quickly forgotten. This study concludes that memory is structured to prioritize the "deep structure" or meaning of a sentence over its "surface structure" or syntax, suggesting that comprehension relies on extracting and retaining the core meaning of information.

{\color{maroon}
The temporal dynamics of semantic versus syntactic memory retention revealed by Sachs demonstrate fundamental principles of human cognition that directly support semantic similarity approaches in educational assessment. After 80-160 syllables of interpolated material, recognition for syntactic changes dropped to near chance levels while semantic recognition remained robust, indicating that individuals naturally retain the meaning or "gist" of information rather than its exact syntactic form.

This natural tendency toward semantic retention over syntactic precision has profound implications for educational technology design. Traditional evaluation methods that rely on verbatim reproduction fail to account for this fundamental cognitive process, potentially penalizing students who demonstrate genuine understanding through paraphrasing and conceptual restructuring. The findings support assessment approaches that recognize semantic equivalence while maintaining academic rigor, providing theoretical justification for similarity-based evaluation systems that align with natural memory processes.
}

\subsection{Cognitive Processes in Memory Retention}

Rashid and Rigas (2006) provide an analysis of memory stages and their relationship to learning and information retention. The study outlines the progression of information through sensory memory, short-term memory, and long-term memory, highlighting the necessity of active cognitive processing for effective retention. Encoding, which involves the initial capture of relevant information, and reviewing, which reinforces the transfer of information to long-term memory, are identified as critical phases in the learning process. This study also discusses retrieval failure as a cause of forgetting, where proactive interference and retroactive interference disrupt recall. 

Proactive interference occurs when previously learned information impairs the learning of new material, while retroactive interference arises when newly acquired knowledge interferes with the recall of older information. These findings emphasize that memory retention is optimized when learners engage in deep processing activities, such as organizing and summarizing information, rather than relying on surface-level strategies. This perspective aligns with the findings of Sachs (1967), reinforcing the importance of meaningful encoding and the prioritization of core content over superficial details.

\subsection{Memory and Recall}

Naim et al. (2020) identified that recall performance grows sublinearly as the list length increases. This sublinear growth highlights a saturation effect, where the number of items retrieved plateaus despite an increase in encoded material. The study also revealed that encoding quality plays a crucial role in determining recall performance. Slower presentation rates allowed participants to encode a higher number of words into memory, which improved the total recall capacity. These findings provide evidence that memory recall adheres to universal laws, offering a structured and deterministic perspective on retrieval processes.

\subsubsection{Retrieval-Based Testing}

Karpicke and Roediger (2008) emphasize that testing is not merely a means of assessment but a potent strategy for enhancing long-term memory retention. Their study debunks the conventional belief that once information is initially learned, further practice adds little value. Instead, repeated retrieval—through testing—promotes deeper encoding and strengthens memory consolidation, highlighting retrieval as an active mechanism essential for sustained learning performance.

\subsubsection{Optimizing Retrieval Practice for Long-term Retention}

Karpicke and Roediger (2008) demonstrated that retrieval practice plays a central role in consolidating memories and improving long-term retention. Their experiments revealed that repeated retrieval significantly outperforms repeated study in enhancing memory performance. This finding validates the importance of active recall over passive encoding strategies.

{\color{maroon}
The systematic investigation by Karpicke and Roediger reveals critical distinctions between study repetition and test repetition effects that challenge conventional educational assumptions. Once students correctly produce learning items, continued studying provides no measurable retention benefit, whereas continued testing generates substantial positive effects on long-term memory consolidation. This differential impact suggests that retrieval practice engages cognitive mechanisms distinct from those involved in passive review.

Furthermore, students' predictions of their own learning performance show no correlation with actual retention outcomes, indicating poor metacognitive awareness of effective learning strategies. This metacognitive blindness suggests that students may persist with ineffective study methods while avoiding retrieval practice that would substantially improve their learning outcomes. These findings provide empirical support for educational systems that prioritize active recall over passive content review.

\subsection{Mathematical Models of Memory Recall}

Naim et al. (2020) demonstrate that human memory recall, despite appearing fragile and unpredictable, follows universal mathematical laws that can be derived from first principles. Their research on free recall paradigms reveals that memory operates as a deterministic process rather than a random phenomenon, with recall performance following predictable mathematical relationships.

The fundamental law of memory recall establishes that the average number of items recalled (R) from M items in memory follows the relationship $R = \sqrt{\frac{3\pi M}{2}}$. This parameter-free prediction emerges from modeling memory as associative search processes operating on random graphs defined by memory representation structure.

The theoretical foundation rests on two principles: memory items are represented by overlapping random sparse neuronal ensembles in dedicated memory networks, and recall follows associative chains where the next item recalled has the largest overlap with the current item. This associative retrieval mechanism provides mathematical validation for similarity-based recall processes, supporting computational approaches to memory assessment that rely on semantic similarity rather than exact matching.
}

\section{Note-Taking Techniques and Tools}

Note-taking serves as an essential tool for facilitating memory retention and comprehension. Various techniques and tools have been developed to assist learners in organizing and processing information effectively.

\subsection{Traditional Note-Taking Techniques}

Rashid and Rigas (2006) examine the comparative effectiveness of traditional note-taking methods, including the Cornell Method, the Outlining Method, and the Mapping Method. These methods provide learners with strategies for organizing content to improve comprehension and retention.

\subsection{Digital Note-Taking and E-Learning Tools}

Rashid and Rigas (2006) further explore the role of technology in enhancing note-taking and memory retention within e-learning environments. The integration of tools such as Microsoft OneNote and Tablet PCs has provided learners with new ways to engage with content, enabling multimedia annotations and portable accessibility.

\subsubsection{Quizlet}

Tran and Nguyen (2022) demonstrate that Quizlet integrates notes, spaced repetition, and active learning into an engaging platform, effectively improving vocabulary retention and long-term memory consolidation.

\subsubsection{ChatGPT and the Generation of Educational Content}

Mishra and Anthony (2023) highlight both the benefits and drawbacks of AI tools like ChatGPT in education, emphasizing the need for responsible integration to enhance critical thinking and active learning processes.

\subsubsection{Self-Regulated Learning by Using Study Applications}

Baars, Khare, and Ridderstap (2022) demonstrate that study applications with well-designed user interfaces support self-regulated learning by enabling goal-setting, time management, and feedback during study sessions.

\section{NLP and Semantic Similarity}

The development of NLP techniques has significantly advanced models that analyze and quantify semantic similarity.

{\color{maroon}
\subsection{Word Embedding Foundations}

The evolution of neural word embeddings represents a foundational shift in computational linguistics, establishing the vector space principles that underpin contemporary sentence embedding methodologies. Church (2017) demonstrates that Word2Vec's widespread adoption stems not from theoretical complexity, but from its practical accessibility and the availability of downloadable implementations that enable immediate application. The method's core innovation lies in representing words as dense vectors in continuous space, where semantic relationships are preserved through geometric relationships rather than symbolic representations.

Mikolov et al. (2013) introduced the skip-gram architecture that learns word representations by predicting contextual words, establishing contrastive learning principles that would later be adapted for sentence-level embeddings. The famous analogy \textit{man : woman :: king : queen} illustrates how vector arithmetic captures semantic relationships, with the mathematical operation $\vec{king} - \vec{man} + \vec{woman} \approx \vec{queen}$ demonstrating the foundation of semantic similarity computation through vector mathematics.

Jatnika, Bijaksana, and Suryani (2019) provide empirical validation of Word2Vec's semantic similarity capabilities through extensive evaluation on WordSim-353 and SimLex-999 datasets. Their analysis reveals that cosine similarity, computed as the angular distance between word vectors, effectively quantifies semantic relationships with correlation scores reaching 0.665 on WordSim-353. This establishes the mathematical foundation for similarity measurement that extends to sentence-level embeddings, where the principle that semantically similar entities should be positioned closer in vector space forms the theoretical basis for contrastive learning objectives used in SimCSE and similar frameworks.
}

\subsection{Text Semantic Similarity}

Corley et al. (2005) provided foundational insights into computational methods for evaluating semantic similarity, focusing on techniques that measure the closeness of meaning between texts.

\subsection{The BERT Framework}

Devlin et al. (2019) emphasize that BERT leverages bidirectional representations of language to capture semantic relationships through its transformer-based architecture.

\subsubsection{SimCSE}

Gao, Yao, and Chen (2021) developed SimCSE, a fine-tuning approach that enhances the quality of sentence embeddings, making them suitable for tasks like semantic similarity evaluation and textual entailment.

{\color{maroon}
SimCSE addresses a fundamental challenge in sentence embeddings known as representation collapse, where different sentences produce nearly identical vector representations, rendering similarity measurements meaningless. Traditional BERT embeddings suffer from this phenomenon due to anisotropic distribution of representations in vector space, where embeddings cluster in narrow regions rather than utilizing the full dimensional space effectively.

The theoretical contribution of SimCSE lies in transforming anisotropic embedding distributions into more isotropic ones. Anisotropic embeddings concentrate in narrow cones within the vector space, limiting their ability to capture fine-grained semantic distinctions necessary for educational assessment. The contrastive learning framework regularizes the embedding space to achieve uniform distribution, where sentences with different meanings occupy distinct regions while semantically similar sentences remain proximate.

Unlike previous approaches requiring complex data augmentation techniques such as back-translation, paraphrasing, or syntactic transformations, SimCSE achieves superior performance using only dropout noise as minimal data augmentation. This approach preserves semantic content while introducing controlled perturbations necessary for contrastive learning. When the same sentence passes through the encoder twice with different dropout masks, the resulting embeddings form natural positive pairs without external modification.

The mathematical foundation of SimCSE rests on achieving uniform distribution property in the embedding space. This property ensures that the similarity function can effectively discriminate between sentences with varying degrees of semantic overlap, eliminating the bias present in anisotropic embeddings where high similarity scores may result from distributional artifacts rather than genuine semantic relationships.
}

\subsubsection{Contrastive Learning}

SimCSE's unsupervised framework generates "positive pairs" by applying independent dropout masks to the same input sentence. Specifically, when a sentence is passed through the BERT encoder twice, the randomness introduced by dropout produces two slightly different embeddings.

{\color{maroon}
SimCSE's contrastive learning framework builds upon foundational principles established in Word2Vec's skip-gram model while extending them to sentence-level representations. The approach creates positive pairs through minimal data augmentation and efficiently generates negative samples through in-batch negatives, eliminating the computational overhead of explicitly creating negative pairs while ensuring diverse negative samples.

For a given sentence $s_i$, all other sentences $s_j$ where $j \neq i$ in the batch serve as negatives, creating rich contrastive signals without additional data preprocessing. This in-batch negative sampling strategy demonstrates continuity with Word2Vec's contrastive learning principles, where the fundamental approach of learning through contrast remains consistent across word-level and sentence-level embeddings.
}

\begin{figure}[ht!]
    \centering
    \[
    \ell_i = - \log \frac{\exp(\text{sim}(h_i, h_i^+)/\tau)}{\sum_{j=1}^N \exp(\text{sim}(h_i, h_j)/\tau)}
    \]
    \caption[Formula for Contrastive Learning]{Formula for Contrastive Learning}
    \label{fig:contrastive_learning}
\end{figure}

\subsubsection{Cosine Similarity}

Cosine similarity quantifies the angular distance between embeddings, ensuring that semantically similar sentences are pulled closer together, while dissimilar ones are pushed apart. The temperature parameter $\tau$ sharpens the similarity distribution, controlling the influence of positive and negative pairs during training.

{\color{maroon}
The temperature parameter $\tau$ in the contrastive learning formula plays a critical role in controlling the sharpness of the similarity distribution. Lower temperature values create more peaked distributions, emphasizing high-confidence similarities, while higher values produce smoother distributions. The optimal temperature balances discriminative power with stability, preventing over-confident predictions that could lead to poor generalization in educational assessment contexts.
}

\begin{figure}[ht!]
    \centering
    \[
    \text{sim}(h_1, h_2) = \frac{h_1^\top h_2}{\|h_1\|\|h_2\|}
    \]
    \caption[Formula for Cosine Similarity]{Formula for Cosine Similarity}
    \label{fig:cosine_similarity}
\end{figure}



\section{Progress Tracking}

Progress tracking refers to systems that enable users to monitor and evaluate their advancement toward specific goals.

\subsection{Progress Tracking as Part of Gameful Design}

Deterding et al. (2011) emphasize the importance of structured systems that guide users through progressive goals and stages of development, helping maintain motivation and engagement.

\subsubsection{The Three-Star Scoring System}

Gastón and Cooper (2017) explore the efficacy of the three-star scoring system, highlighting its ability to improve user engagement and performance in educational and gamified applications.

{\color{maroon}
\section{Web Development Framework Literature}

The selection of appropriate web development technologies is crucial for building scalable, maintainable, and user-friendly educational applications. This section reviews the literature surrounding modern web development frameworks and libraries that form the technical foundation of contemporary educational platforms.

\subsection{Next.js and React Framework}

Vercel (2023) emphasizes that Next.js provides a production-ready React framework with built-in optimizations, server-side rendering capabilities, and automatic code splitting that enhances application performance and user experience. The framework's hybrid static and server rendering approach enables developers to create fast, SEO-optimized applications suitable for educational platforms requiring both dynamic content and static resources.

React's component-based architecture, as documented by Facebook (2023), facilitates the development of reusable UI components that can be efficiently maintained and scaled. The library's virtual DOM implementation optimizes rendering performance, which is particularly beneficial for educational applications that frequently update content and user interfaces.

\subsection{TypeScript Integration}

Microsoft (2023) demonstrates that TypeScript's static type checking reduces runtime errors and improves code maintainability in large-scale applications. The type safety provided by TypeScript is particularly valuable in educational software development, where data integrity and reliable user interactions are paramount.

Bierman et al. (2014) highlight TypeScript's gradual typing system, which allows developers to incrementally add type annotations to existing JavaScript codebases, facilitating smoother migration processes and reducing development risks.

\subsection{Data Visualization Technologies}

Chart.js, as documented by Chart.js Community (2023), provides a flexible and responsive charting library that enables the creation of interactive data visualizations. In educational applications, effective data visualization supports learning analytics by presenting user progress, performance metrics, and learning patterns in comprehensible formats.

Bostock et al. (2011) established that interactive visualizations enhance user engagement and comprehension, particularly in educational contexts where learners benefit from immediate visual feedback on their progress and performance.

\subsection{Styling and User Interface}

Wind et al. (2023) document Tailwind CSS as a utility-first CSS framework that accelerates development while maintaining design consistency. The framework's approach to styling through utility classes enables rapid prototyping and maintains design system coherence across educational applications.

Adams (2021) argues that consistent and responsive user interfaces are critical for educational applications, as they reduce cognitive load and allow learners to focus on content rather than navigation challenges.

\subsection{Database and Backend Technologies}

PostgreSQL's robust relational database capabilities, as documented by PostgreSQL Global Development Group (2023), provide the reliability and ACID compliance necessary for educational applications that manage user data, learning progress, and assessment results. The database's support for complex queries and data integrity constraints ensures accurate tracking of educational metrics.

Cloudinary's cloud-based media management platform, according to Cloudinary (2023), enables efficient handling of educational multimedia content, including automatic optimization and responsive delivery that adapts to various devices and network conditions commonly encountered in educational settings.

}

\section{Summary}

The studies reviewed in this chapter highlight key advancements in note-taking strategies, memory processes, NLP, and progress tracking. These insights establish a foundation for the proposed system, integrating effective learning strategies and semantic similarity evaluation to optimize recall and knowledge retention.

The table compares various studies based on their focus on specific aspects
or features that are relevant to the proposed approach. It highlights which aspects are
addressed and identifies gaps that the proposed approach aims to fill by integrating
these features into a unified framework.

\begin{figure}[ht!]
    \centering
    \includegraphics[width=0.9\textwidth]{comparison_table.png} % Replace with your image file's name
    \caption[Comparison Table of Related Literature]{Comparison Table of Related Literature}
    \label{fig:comparison_table}
\end{figure}

